\scp{\tsc{Rote-Meto}}{04-02}
The \tsc{Rote-Meto} node contains the languages of Rote Island
as well as the \tsc{Meto} language/dialect cluster which is spoken
across the western part of the island of Timor,
including the enclave of Oecusse (politically part of Timor-Leste).
The sound changes and history of the \tsc{Rote-Meto} languages
have been explored in detail in a series of publications
\citep{Edwards2016,Edwards2018a,Edwards2018c}
and over 1,000 Proto-Rote-Meto reconstructions
have been made \citep{Edwards2021}.

Although the languages of Rote have sometimes been treated as
a single language, they have undergone different sound changes
and have differences in grammar,
lexicon, and phonology to such a degree that intelligibility
between non-adjacent varieties is often blocked
(\citealp[47]{GrimesTherikGrimesJacob1997}, \citealp[233]{Fox2016}, \citealp[2--4]{Tamelan2021}).
Indeed, there is not even good phonological or morphological
evidence for treating all the languages of Rote as unitary node
below the level of \tsc{Rote-Meto}.
Similarly the Meto language/dialect cluster has sometimes been
treated as a single language (labelled Dawan[ese],
Timorese, or Atoni) and there is a remarkable level of lexical
cohesion among varieties of Meto,
with cognacy between the most widely separated varieties (Kopas
and Kusa-Manea) ranging between 77--83{\%}.
Nonetheless, there is not insignificant diversity among non-basic
lexical items, phonology, semantics, morphology, and morphosyntax
and speakers of different varieties frequently report difficulty
communicating with one another \citep[6]{Edwards2020}.

There are two sound changes which define the \tsc{Rote-Meto}
node within \tsc{Timor-Babar}: *wa > \it{o} in a lexically specific
set of words, and the lowering of many high vowels before final *R,
after which *R > {\0}.
Of the fifteen examples of PMP *wa so far attested in \tsc{Rote-Meto},
nine show *wa > \it{o}.\footnote{
		The two words not shown in \trf{04-03} which
		also have *wa > \it{o} are PMP *hawak > Proto-Rote-Meto **ao-k
		‘body’ and PWMP *qambawaŋ > Proto-Rote-Meto **mbao ‘mango’.}
Seven of these examples are given in \trf{04-03}.\footnote{
		The seven
		words which do not have *wa > \it{o} in \tsc{Rote-Meto} are:
		PMP *walu > Proto-Rote-Meto **falu ‘eight’,
		*waRi > **fai ‘day,
		time’, PWMP *bayawak > **ɓaiafa ‘monitor lizard’,
		*huaji > **waji > **waɗi ‘same sex younger sibling’,
		*wani = **wani ‘bee’,
		*kali-wati > **kalati ‘earthworm’,
		and *lawaq > **ɓolau ‘spider’ (source of initial syllable unknown).}

\begin{table}\tcp{\tsc{Rote-Meto} *wa > \it{o}}{04-03}
\begin{adjustbox}{width=\textwidth}
	\begin{threeparttable}\st{2pt}
	\begin{tabular}{llllllll}\lsptoprule
local gloss	&	{\hr}water	&	{\hr}root	&	{\hr}dry in sun	&	{\hr}right	&	{\hr}nine	&	{\hr}marry	&	{\hr}viper\\
PMP	&	{\hr}*\tbr{wa}hiyəR	&	{\hr}*\tbr{wa}kaR\su{‡}	&	{\hr}*pa-\tbr{wa}Ri	&	{\hr}*ka\tbr{wa}nan	&	{\hr}*si\tbr{wa}	&	{\hr}*qasa\tbr{wa}	&	{\hr}*sa\tbr{wa}\\
Proto-R-M	&	**\tbr{o}e	&	**\tbr{o}ka-k	&	**h\tbr{o}i	&	**k\tbr{o}na	&	**si\tbr{o}	&	**sa\tbr{o}	&	**kaisa\tbr{o}\su{Δ}\\ \midrule
Rikou	&	{\hr\hr}\tbr{o}e	&	{\hr\hr}\tbr{o}ka-ʔ	&	{\hr\hr}h\tbr{o}i	&	{\hr\hr}ʔ\tbr{o}na	&	{\hr\hr}si\tbr{o}	&	{\hr\hr}sa\tbr{o}	&	\\
Bilbaa	&	{\hr\hr}\tbr{o}e	&	{\hr\hr}\tbr{o}ka-ʔ	&	{\hr\hr}h\tbr{o}i	&	{\hr\hr}k\tbr{o}na	&	{\hr\hr}si\tbr{o}	&	{\hr\hr}sa\tbr{o}	&	\\
Termanu	&	{\hr\hr}\tbr{o}e	&	{\hr\hr}\tbr{o}ka-k	&	{\hr\hr}h\tbr{o}i	&	{\hr\hr}k\tbr{o}na	&	{\hr\hr}si\tbr{o}	&	{\hr\hr}sa\tbr{o}	&	{\hr\hr}kaisa\tbr{o}\\
Bokai	&	{\hr\hr}\tbr{o}e	&	{\hr\hr}\tbr{o}ka-k	&	{\hr\hr}h\tbr{o}i	&	{\hr\hr}k\tbr{o}na	&	{\hr\hr}si\tbr{o}	&	{\hr\hr}sa\tbr{o}	&	{\hr\hr}kaisa\tbr{o}\\
Ba{\Q}a, Tii	&	{\hr\hr}\tbr{o}e	&	{\hr\hr}\tbr{o}ka-k	&	{\hr\hr}h\tbr{o}i	&	{\hr\hr}k\tbr{o}na	&	{\hr\hr}si\tbr{o}	&	{\hr\hr}sa\tbr{o}	&	{\hr\hr}kaisa\tbr{o}\\
Dengga	&	{\hr\hr}\tbr{o}e	&	{\hr\hr}ʔ\tbr{o}ka-ʔ	&	{\hr\hr}h\tbr{o}i	&	{\hr\hr}ʔ\tbr{o}na	&	{\hr\hr}si\tbr{o}	&	{\hr\hr}sa\tbr{o}	&	{\hr\hr}aisa\tbr{o}\\
Dela-O.	&	{\hr\hr}\tbr{o}e	&	{\hr\hr}ʔ\tbr{o}ka-ʔ	&	{\hr\hr}h\tbr{o}i	&	{\hr\hr}ʔ\tbr{o}na	&	{\hr\hr}si\tbr{o}	&	{\hr\hr}sa\tbr{o}	&	{\hr\hr}aisa\tbr{o}\\
Roi{\Q}is	&	{\hr\hr}\tbr{o}e	&	\hr\cg{(baʔa-n)}	&	{\hr\hr}n-h\tbr{o}e	&	\hr\cg{(ʔneʔu)}	&	{\hr\hr}se\tbr{o}	&	{\hr\hr}n-sa\tbr{o}	&	{\hr\hr}sa\tbr{o}\\
Kotos	&	{\hr\hr}\tbr{o}e	&	\hr\cg{(ʔbaʔa-f)}	&	{\hr\hr}n-h\tbr{o}e	&	\hr\cg{(ʔneʔu)}	&	{\hr\hr}se\tbr{o}	&	{\hr\hr}n-sa\tbr{o}	&	{\hr\hr}ʔsa\tbr{o}\\
Molo	&	{\hr\hr}\tbr{o}e	&	\hr\cg{(ʔbaʔat)}	&	{\hr\hr}n-h\tbr{o}i	&	\hr\cg{(ʔneʔu)}	&	{\hr\hr}se\tbr{o},	&	{\hr\hr}n-sa\tbr{o}	&	{\hr\hr}uʔsa\tbr{o}, \\
	&		&		&		&		&	 {\hr\hr}si\tbr{o}	&	&	{\hr\hr}aʔsa\tbr{o}\\
Kusa-M.	&	{\hr\hr}\tbr{o}a\su{†}	&	\hr\cg{(baʔa-f)}	&	{\hr\hr}h\tbr{o}i	&	\hr\cg{(neʔu)}	&	{\hr\hr}si\tbr{o}	&		&	{\hr\hr}sa\tbr{o}\\
		\lspbottomrule
	\end{tabular}
		\begin{tablenotes}\tnsize
			\item[†]	{Kusa-Manea data is from the Upper Manulea dialect in which
								**e > \it{a} /V{\gap}{\#} is a regular change.}
			\item[‡]	{The \tsc{Meto} words meaning ‘root’ appear to be connected
								with Funai Helong \it{kbakat} ‘root’.
								This word and the \tsc{Meto} words may further be irregularly
								connected with PMP *wakaR.}
			\item[Δ]	{The initial element of Proto-Rote-Meto **kaisao ‘viper’
								may be from PMP *kayu ‘tree, wood’.}
		\end{tablenotes}
	\end{threeparttable}
\end{adjustbox}
\end{table}

Although some other \tsc{Timor-Babar} languages also have *wa
> \it{o} in some words,
there is no reason to suppose that this is due to a period of
common development shared between \tsc{Rote-Meto} and these languages.
Thus, for instance, Kisar also has *wa > \it{o} in *wahiyəR >
\it{oir} ‘water’ and *qasawa > \it{hoo} ‘marry’.
However, other cognates in \tsc{Southwest Maluku} show
that this sound change occurred in Kisar after the break-up of Proto-Southwest Maluku.
It thus cannot be assigned to a single common ancestor of \tsc{Rote-Meto} and \tsc{Southwest Maluku}.

Likewise, \tsc{Wetar-Atauro} languages have word medial *wa >
\it{o} in some reflexes (see \trf{04-18}).
However, word-initial *w is lost (i.e.
*w > {\0}) or is retained without any change to the following vowel.
Word medial *wa > \it{o} is thus best viewed as parallel development.

Helong has *wa > \it{u} in *wahiyəR > \it{ui} ‘water’
and *pa-waRi > \it{huiʔ} ‘dry in sun’.
Helong does not allow sequences of a mid-vowel followed by a
high vowel, and thus both these words could be from intermediate **oi
and **hoiʔ respectively, though they could also be from intermediate **ue or **hueʔ.
Thus, in these two words Helong may share *wa > **o with \tsc{Rote-Meto}.
However, *pa-waRi > \it{huiʔ} ‘dry in sun’ has irregular *R >
{\0} instead of regular *R > \it{l}
and is thus probably a borrowing from a \tsc{Rote-Meto} language.
Thus, at best, there is one word,
*wahiyəR ‘water’, for which \tsc{Rote-Meto}
and Helong may share *wa > **o.

An additional sound change which provides evidence that \tsc{Rote-Meto}
is a valid node within \tsc{Timor-Babar} is sporadic lowering
of high vowels to mid before *R.
Six examples of this change are shown in \trf{04-04}. PMP *R was
subsequently lost in most \tsc{Rote-Meto} words.

\begin{table}\tcp{\tsc{Rote-Meto} V\tsc{[+high]} > V\tsc{[+mid]} /{\gap}R{\#}}{04-04}
%\begin{adjustbox}{width=\textwidth}
	\begin{threeparttable}\st{2.5pt}
	\begin{tabular}{lllllll}\lsptoprule
local gloss	&	{\hr}tail	&	{\hr}egg	&	{\hr}calcium	&	{\hr}grain head	&	{\hr}way\su{‡}	&	{\hr}hearth\\
PMP	&	{\hr}*ik\tbr{uR}	&	{\hr}*qatəl\tbr{uR}	&	{\hr}*qap\tbr{uR}	&	{\hr}*bul\tbr{iR}	&	{\hr}*qən\tbr{uR}	&	{\hr}*dap\tbr{uR}\\
Proto-R-M	&	**ik\tbr{o}-k	&	**təl\tbr{o}-k	&	**ah\tbr{o}	&	**mbul\tbr{e}-k	&	**eno-k	&	**raho\su{ǁ}\\ \midrule
Rikou	&	{\hr\hr}ik\tbr{o}-ʔ	&	{\hr\hr}tol\tbr{o}-ʔ	&	{\hr\hr}a\tbr{o}	&	{\hr\hr}pul\tbr{e}-ʔ	&	{\hr\hr}en\tbr{o}-ʔ	&	{\hr\hr}ra\tbr{o}\\
Bilbaa	&	{\hr\hr}ik\tbr{o}-ʔ	&	{\hr\hr}tol\tbr{o}-ʔ	&	{\hr\hr}a\tbr{o}	&	{\hr\hr}pul\tbr{e}-ʔ	&	{\hr\hr}en\tbr{o}-ʔ	&	{\hr\hr}la\tbr{o}\\
Termanu	&	{\hr\hr}ik\tbr{o}-k	&	{\hr\hr}tol\tbr{o}-k	&	{\hr\hr}haʔ\tbr{o}	&	{\hr\hr}pul\tbr{e}-k	&	{\hr\hr}en\tbr{o}-k	&	{\hr\hr}laʔ\tbr{o}\\
Bokai	&	{\hr\hr}ik\tbr{o}-k	&	{\hr\hr}tol\tbr{o}-k	&	{\hr\hr}haʔ\tbr{o}	&	{\hr\hr}pul\tbr{e}-k	&	{\hr\hr}en\tbr{o}-k	&	{\hr\hr}laʔ\tbr{o}\\
Ba{\Q}a	&	{\hr\hr}ik\tbr{o}-k	&	{\hr\hr}tol\tbr{o}-k	&	{\hr\hr}haʔ\tbr{o}	&	{\hr\hr}mpul\tbr{e}-k	&		&	{\hr\hr}laʔ\tbr{o}\\
Tii	&	{\hr\hr}ik\tbr{o}-k	&	{\hr\hr}tol\tbr{o}-k	&	{\hr\hr}aʔ\tbr{o}	&	{\hr\hr}mbul\tbr{e}-k	&	{\hr\hr}en\tbr{o}-k	&	{\hr\hr}raʔ\tbr{o}\\
Dengga	&	{\hr\hr}ik\tbr{o}-ʔ	&	{\hr\hr}tel\tbr{o}-ʔ	&	{\hr\hr}a\tbr{o}	&	{\hr\hr}mbul\tbr{e}-ʔ	&		&	\\
Dela-Oenale	&	{\hr\hr}ik\tbr{o}-ʔ	&	{\hr\hr}tel\tbr{o}-ʔ	&	{\hr\hr}a\tbr{o}	&	{\hr\hr}mbul\tbr{e}-ʔ	&		&	\\
Roi{\Q}is	&	{\hr\hr}i͡uk\tbr{u}-f \su{†}	&	{\hr\hr}ten\tbr{o}ʔ	&	{\hr\hr}a\tbr{o}	&	{\hr\hr}puniʔ\su{†}	&	{\hr\hr}en\tbr{o}ʔ	&	\\
Kotos	&	{\hr\hr}ik\tbr{u}-f \su{†}	&	{\hr\hr}ten\tbr{o}ʔ, 	&	{\hr\hr}a\tbr{o}	&	{\hr\hr}pun\tbr{e}ʔ,	&	{\hr\hr}en\tbr{o}ʔ	&	{\hr\hr}auf na\tbr{o}\su{Δ}\\
	&	 	&	{\hr\hr}tek\tbr{o}ʔ	&		&	{\hr\hr}pun\tbr{i}ʔ	&		&	\\
Molo	&	{\hr\hr}ik\tbr{o}-f	&	{\hr\hr}tek\tbr{o}ʔ	&	{\hr\hr}a\tbr{o}	&	{\hr\hr}pun\tbr{e}ʔ	&	{\hr\hr}en\tbr{o}ʔ	&	\\
Kusa-Manea	&	{\hr\hr}ik\tbr{u}-f \su{†}	&	{\hr\hr}ten\tbr{o}ʔ	&	{\hr\hr}a\tbr{o}	&		&		&	\\
		\lspbottomrule
	\end{tabular}
		\begin{tablenotes}\tnsize
			\item[†]	{Kotos Amarasi, Roi{\Q}is Amarasi,
								and Kusa-Manea have sporadic raising of final mid-vowels to high
								after a high vowel.}
			\item[‡]	{The reflexes of *qənuR in Rote mean ‘way’,
								the reflexes in \tsc{Meto} mean ‘door, opening, gate’.}
			\item[Δ]	{Kotos Amarasi \it{auf nao} is ‘ash in the hearth’.}
			\item[{ǁ}]	{PMP *dapuR > Proto-Rote-Meto **raho has minority
								*d > **r \citep[313]{Edwards2021}.}
		\end{tablenotes}
	\end{threeparttable}
%\end{adjustbox}
\end{table}

Within \tsc{Timor-Babar}, Lakalei and Idate (\srf{04-04}) also
have vowel lowering before *R in reflexes of *qatəluR ‘egg’ (Lakalei,
Idate \it{manu telor} ‘egg’)
and *ikuR ‘tail’ (Lakalei \it{hiʔon},
Idate \it{yiʔan} ‘tail’) but do not have it in reflexes of *qapuR
‘calcium’ (Lakalei \it{aur}, Idate \it{waʔur} ‘calcium’).
No other \tsc{Timor-Babar} languages are known to have lowering
of final high vowels,
with the possible exception of Helong *qapuR > \it{ao} ‘calcium’
which is identical to Meto \it{ao} ‘calcium’ from which it may be a loan.
While we could subgroup \tsc{Lakalei-Idate}
and \tsc{Rote-Meto} together on the basis of lowering of high
vowels to mid, this would currently rest on only two words.

All languages of the \tsc{Central Timor} subgroup (\crf{ch:CT})
also have lowering of high vowels to mid before final *R,
with the exception of Welaun.
Vowel lowering before *R is thus a regional feature of west and central Timor.
Nonetheless, even if vowel lowering before *R in \tsc{Rote-Meto},
\tsc{Lakalei-Idate}, and \tsc{Central Timor} languages is due
to contact, the fact that it had occurred in a lexically specific set of
words by the time of Proto-Rote-Meto can be taken as supporting
evidence for \tsc{Rote-Meto} as a subgroup.
See \citet[82--86]{Edwards2021} for more detailed discussion
of sound changes which allow us to identify \tsc{Rote-Meto} as
comprising a single node.

One morphological feature which may provide additional evidence
for \tsc{Rote-Meto} is the innovation of a reciprocal prefix **ma-.
Illustrated in \qf{04-04}--\qf{04-06} below.
We have yet to encounter a similar **ma- \tsc{recip} prefix in surrounding neighbours.
However, this prefix may be a retention of *maR-
for which a reciprocal function has been reconstructed \citep{Bril2019}.

{\st{3pt}\wg\begin{xltabular}[l]{\linewidth}{>{\hspace{-\tabcolsep}}l<{\hspace{-\tabcolsep}}>{\hspace{-\tabcolsep}\enspace}llll}
\tbx{04-04}	&\mc{3}{l}{Rikou (eastern Rote) \it{ma-} \tsc{recip}}& \mc{1}{r}{\citep{NakoBanamtuan2014}}\\
						&\it{resi}	&‘argue’&\it{\tbr{ma}-re{\tl}resi-ʔ}	&‘argue with each other, debate’\\
						&\it{etu}	&‘separate’&\it{ra-\tbr{ma}-etu-ʔ}	&‘divorce one another’\\
\tbx{04-05}	&\mc{3}{l}{Dela (western Rote) \it{ma-} \tsc{recip}}& \mc{1}{r}{\citep[156]{Tamelan2021}}\\
						&\it{helu}	&‘promise’&\it{\tbr{ma}-helu-ʔ}	&‘promise each other’\\
						&\it{hela}	&‘separate, divorce’&\it{\tbr{ma}-hela-ʔ}	&‘divorce one another’\\
						&\it{ʔira}	&‘convince, forbid’&\it{\tbr{ma}-ʔira-ʔ}	&‘agree with each other’\\
						&\it{toʔu}	&‘hold’&\it{\tbr{ma}-toʔu-ʔ}	&‘hold each other’\\
						&\it{tuŋga}	&‘follow’&\it{\tbr{ma}-tu{\tl}tuŋga-ʔ}	&‘follow each oher’\\
\tbx{04-06}	&\mc{3}{l}{Amanuban (\tsc{Meto}) \it{ma-} \tsc{recip}}& \mc{1}{r}{(Grimes fieldnotes)}\\
						&\it{akol}	&‘plot’&\it{n-\tbr{ma}-ʔakol=iin}	&‘plot together’\\
						&\it{naka}	&‘border’&\it{n-\tbr{ma}-naka=n}	&‘share a common border’\\
						&\it{nena}	&‘listen, hear’&\it{n-\tbr{ma}-nena=n}	&‘listen to one another’\\
						&\it{tabe}	&‘greet’&\it{n-\tbr{ma}-tabe=n}	&‘greet each other’\\
						&\it{toe}		&‘argue’&\it{n-\tbr{ma}-toe=n}	&‘argue with each other’\\
\end{xltabular}\addtocounter{table}{-1}}

Within \tsc{Rote-Meto} we can identify a \tsc{West Rote-Meto}
node containing Dela-Oenale, Dengga, and \tsc{Meto}, as well as a \tsc{Nuclear Rote} node
containing all other Rote languages.
\tsc{West Rote-Meto} is defined by six changes given in (7)
and \tsc{Nuclear Rote} by three changes given in (8).\footnote{
		\tsc{West
		Rote-Meto} and \tsc{Nuclear Rote} also have a
		certain amount of exclusively shared vocabulary.
		In his reconstruction of Proto-Rote Meto \citet[70]{Edwards2021}
		has 37 Proto-West Rote-Meto reconstructions
		and 23 Proto-Nuclear Rote reconstructions.
		While some of these are retentions which have been lost in the other branch,
		others may be shared innovations.}
 Examples of Proto-Rote-Meto**d,
**wa, and final-syllable **ə are given in \trf{04-05}.\footnote{
		Proto-Rote-Meto
		**d, **ɗ, **r, and **ʤ each have multiple sources due to unconditioned
		splits in PMP proto-phonemes.
		The sound changes from PMP yielding these proto-phonemes are:
		*d > **d, **ɗ, **r /{\#}{\gap}; **d > **ɗ,
		**r /V{\gap}V; *z > **ɗ,
		**nd, **d; *j > **ɗ,
		**d, **ʤ, **r; and *R/*r > {\0}, **r.}

{\wg\tblr{>{\hspace{-\tabcolsep}}l<{\hspace{-\tabcolsep}}>{\hspace{-\tabcolsep}\enspace}lllll}
\tbx{04-07}	&	\mc{5}{l}{\tsc{West Rote-Meto} innovations}\\
						&a.	&**ə			&>&\it{a}&/σ\und{σ}\#\\
						&b.	&**wa			&>&\it{o}&/\\
						&c.	&**d/**ʤ	&>&**r&/\\
						&d.	&**b			&>&\it{f}&/{\#}{\gap}\\
						&e.	&**a			&>&\it{a, e}&/{\gap}{\#}\\
						&f.	&**k			&>&\it{h, k}&/{\#}{\gap}\\
\tbx{04-08}	&	\mc{5}{l}{\tsc{Nuclear Rote} innovations}\\
						&a.	&**ə						&>&\it{e}&/σ\und{σ}\#\\
						&b.	&**wa						&>&\it{fa}&/\\
						&c.	&**d/**ɗ/**ʤ	&>&\it{ɗ {\tl} d}&/\\
\end{tabular}\mbox{}\\}

\begin{table}\tcp{\tsc{Rote-Meto} **d, **wa, and final-syllable **ə}{04-05}
%\begin{adjustbox}{width=\textwidth}
	\begin{threeparttable}\st{3pt}
	\begin{tabular}{lllllll}\lsptoprule
local gloss	&	gallbladder	&	bathe	&	inside	&	navel	&	ySib\su{Δ}	&	bee\su{ǁ}\\
PMP	&	{\hr}*qapə\tbr{j}u	&	{\hr}*\tbr{d}iRus	&	{\hr}*\tbr{d}al\tbr{ə}m\su{†}	&	{\hr}*pus\tbr{ə}j	&	{\hr}*h\tbr{ua}ji	&	{\hr}*\tbr{wa}ni\\
Proto-RM	&	**he\tbr{d}u-k	&	**na-\tbr{d}iu	&	**\tbr{d}al\tbr{ə}-k	&	**hus\tbr{ə}-k	&	**\tbr{wa}ɗi	&	**\tbr{wa}ni\\ \midrule
Proto-Nuc. R.	&	**he\tbr{d}u-k	&	**na-\tbr{d}iu	&	**\tbr{d}al\tbr{e}-k	&	**hus\tbr{e}-k	&	**\tbr{fa}di-k	&	**\tbr{fa}ni\\ \midrule
Rikou	&	{\hr\hr}he\tbr{d}u-ʔ	&	{\hr\hr}na-\tbr{d}iu	&	{\hr\hr}\tbr{d}al\tbr{e}-ʔ	&	{\hr\hr}ʔus\tbr{e}-ʔ	&	{\hr\hr}\tbr{fa}di-ʔ	&	{\hr\hr}\tbr{fa}ni\\
Bilbaa	&	{\hr\hr}he\tbr{d}u-ʔ	&	{\hr\hr}na-\tbr{d}iu	&	{\hr\hr}\tbr{d}al\tbr{e}-ʔ	&	{\hr}(ʔ)us\tbr{e}-ʔ	&	{\hr\hr}\tbr{fa}di-ʔ	&	{\hr\hr}\tbr{fa}ni\\
Termanu	&	{\hr\hr}he\tbr{d}u-k	&	{\hr\hr}na-\tbr{d}iu	&	{\hr\hr}\tbr{d}al\tbr{e}-k	&	{\hr}(ʔ)us\tbr{e}-k	&	{\hr\hr}\tbr{fa}di-k	&	{\hr\hr}\tbr{fa}ni\\
Bokai {\&} Ba{\Q}a	&	{\hr\hr}he\tbr{d}u-k	&	{\hr\hr}na-\tbr{d}iu	&	{\hr\hr}\tbr{d}al\tbr{e}-k	&	{\hr}(ʔ)us\tbr{e}-k	&	{\hr\hr}\tbr{fa}di-k	&	{\hr\hr}\tbr{fa}ni\\
Lole {\&} Tii	&	{\hr\hr}he\tbr{ɗ}u-k	&	{\hr\hr}na-\tbr{ɗ}iu	&	{\hr\hr}\tbr{ɗ}al\tbr{e}-k	&	{\hr\hr}hus\tbr{e}-k	&	{\hr\hr}\tbr{fa}ɗi-k	&	{\hr\hr}\tbr{fa}ni\\ \midrule
Proto-WRM	&	**e\tbr{r}u-ʔ	&	**na-\tbr{r}iu	&	**\tbr{r}al\tbr{a}-ʔ	&	**us\tbr{a}-ʔ	&	**\tbr{o}ɗi-ʔ	&	**\tbr{o}ni\\ \midrule
Dela-Oenale	&	{\hr\hr}e\tbr{r}u-ʔ	&	{\hr\hr}na-\tbr{r}iu	&	{\hr\hr}\tbr{r}al\tbr{a}-ʔ	&	{\hr\hr}ʔus\tbr{a}-ʔ	&	{\hr\hr}ʔ\tbr{o}ɗi-ʔ	&	{\hr\hr}\tbr{o}ni\\
Dengga	&	{\hr\hr}e\tbr{l}u-ʔ	&	{\hr\hr}na-\tbr{l}iu	&	{\hr\hr}\tbr{l}al\tbr{a}-ʔ	&	{\hr\hr}ʔus\tbr{a}-ʔ	&	{\hr\hr}ʔ\tbr{o}ɗi-ʔ	&	{\hr\hr}\tbr{o}ni\\
Roi{\Q}is	Amarasi&	{\hr\hr}e͡u\tbr{n}u-f	&	{\hr\hr}na-\tbr{n}iu	&	{\hr\hr}\tbr{n}an\tbr{a}ʔ	&	{\hr\hr}usu-f \su{‡}	&	{\hr\hr}\tbr{o}riʔ	&	{\hr\hr}\tbr{o}ni\\
Kotos Amarasi	&	{\hr\hr}e\tbr{n}u-f	&	{\hr\hr}na-\tbr{n}iu	&	{\hr\hr}\tbr{n}an\tbr{a}-f	&	{\hr\hr}us\tbr{a}-f	&	{\hr\hr}\tbr{o}ri-f	&	{\hr\hr}\tbr{o}ni\\
Molo	&	{\hr\hr}e\tbr{n}uʔ	&	{\hr\hr}na-\tbr{n}iu	&	{\hr\hr}\tbr{n}an\tbr{a}ʔ	&	{\hr\hr}us\tbr{a}ʔ	&	{\hr\hr}\tbr{o}li-f	&	{\hr\hr}\tbr{o}ni\\
Kusa-Manea	&	{\hr\hr}e\tbr{n}u-f	&	{\hr\hr}na-\tbr{n}iu	&	{\hr\hr}\tbr{n}a\tbr{a}n	&	{\hr\hr}us\tbr{a}-f	&	{\hr\hr}\tbr{o}riʔ	&	{\hr\hr}\tbr{o}ni\\
		\lspbottomrule
	\end{tabular}
		\begin{tablenotes}\tnsize
			\item[†]	{Reflexes of *daləm also refer to the seat of emotions in Rote languages (see \srf{02-03-01-06}).
								Kusa-Manea \it{naan} ‘inside’ has final CV → VC metathesis ---
								a productive process in \tsc{Meto}.}
			\item[‡]	{Roi{\Q}is Amarasi \it{usu-f} ‘navel’ has assimilation of final
								/a/ in final closed syllables to the quality of the stressed vowel.
								This is a productive process in Roi{\Q}is Amarasi \citep[181]{Edwards2020}.}
			\item[Δ]	{Reflexes of *huaji mean ‘same sex younger sibling’.}
			\item[{ǁ}]	{Reflexes of *wani mean ‘bee,
								honey, sugar’ in \tsc{Meto}.}
		\end{tablenotes}
	\end{threeparttable}
%\end{adjustbox}
\end{table}

Additional evidence for \tsc{Nuclear Rote} as a valid node comes
from changes in the verb agreement prefixes.
\tsc{Rote-Meto} has two paradigms of verb agreement paradigms:
one with prefixes shaped CV-,
and one with prefixes shaped C-.
The \tsc{Nuclear Rote} languages have all undergone reduction
of the vowel of the CV- prefixes to /a/; that is \tsc{1sg} **ku-
> **ka-, \tsc{2sg} **mu- > **ma-,
and \tsc{1pl.excl}/\tsc{2pl} **mi- > **ma-.\footnote{
		The more
		conservative forms of the prefixes with original vowels are retained
		in \tsc{Nuclear Rote} on reflexes of
		*kaən < **-ha ‘eat’ \citep[165]{Edwards2021}.}
 \tsc{West Rote-Meto} languages retain the original vowels.
CV- verb agreement prefixes in the \tsc{Rote-Meto} languages
are shown in \trf{04-06} along with their likely PMP pronominal
etyma taken from \citet[36]{Ross2002}.

\begin{table}
	\centering\tcp{\tsc{Rote-Meto} CV- verb agreement prefixes}{04-06}
	\st{2pt}\begin{tabular}{llllllll}\lsptoprule
	&	\hr\tsc{1sg}	&	\hr\tsc{2sg}	&	\hr\tsc{3sg}	&	\hr\tsc{1pl.incl}	&	\hr\tsc{1pl.excl}	&	\hr\tsc{2pl}	&	\hr\tsc{3pl}\\
PMP	&	{\hr}*aku	&	{\hr}*kamu	&	{\hr}*ia, {\hr}*nia	&	{\hr}*kita	&	{\hr}*kami	&	{\hr}*miu	&	{\hr}*sida\\
Proto-RM	&	**ku-	&	**mu-	&	**i-, **na-	&	**ta-	&	**mi-	&	**mi-	&	**ra-\\ \midrule
Proto-Nuc. Rote	&	**ka-	&	**ma-	&	**i-, **na-	&	**ta-	&	**ma-	&	**ma-	&	**ra-\\ \midrule
Rikou {\&} Oepao	&	{\hr\hr}ʔa-	&	{\hr\hr}ma-	&	{\hr\hr}na-	&	{\hr\hr}ta-	&	{\hr\hr}ma-	&	{\hr\hr}ma-	&	{\hr\hr}ra-\\
Bilbaa {\&} Bokai	&	{\hr\hr}ka-	&	{\hr\hr}ma-	&	{\hr\hr}na-	&	{\hr\hr}ta-	&	{\hr\hr}ma-	&	{\hr\hr}ma-	&	{\hr\hr}la-\\
Korbafo	&	{\hr\hr}ka-	&	{\hr\hr}ma-	&	{\hr\hr}na-	&	{\hr\hr}ta-	&	{\hr\hr}ma-	&	{\hr\hr}ma-	&	{\hr\hr}la-\\
Termanu {\&} Ba{\Q}a	&	{\hr\hr}ʔa-	&	{\hr\hr}ma-	&	{\hr\hr}na-	&	{\hr\hr}ta-	&	{\hr\hr}ma-	&	{\hr\hr}ma-	&	{\hr\hr}la-\\
Lole	&	{\hr\hr}ʔa-	&	{\hr\hr}ma-	&	{\hr\hr}ni-, i-, na-	&	{\hr\hr}ta-	&	{\hr\hr}ma-	&	{\hr\hr}ma-	&	{\hr\hr}la-\\
Tii	&	{\hr\hr}ʔa-	&	{\hr\hr}ma-	&	{\hr\hr}i-, na-	&	{\hr\hr}ta-	&	{\hr\hr}ma-	&	{\hr\hr}ma-	&	{\hr\hr}ra-\\ \midrule
Proto-WRM	&	**ku-	&	**mu-	&	**na-	&	**ta-	&	**mi-	&	**mi-	&	**ra-\\ \midrule
Dela-Oenale	&	{\hr\hr}ʔu-	&	{\hr\hr}mu-	&	{\hr\hr}na-	&	{\hr\hr}ta-	&	{\hr\hr}mi-	&	{\hr\hr}mi-	&	{\hr\hr}ra-\\
Dengga	&	{\hr\hr}ʔu-	&	{\hr\hr}mu-	&	{\hr\hr}na-	&	{\hr\hr}ta-	&	{\hr\hr}mi-	&	{\hr\hr}mi-	&	{\hr\hr}la-\\
Roi{\Q}is Amarasi	&	{\hr\hr}ku-	&	{\hr\hr}mu-	&	{\hr\hr}na-	&	{\hr\hr}ta-	&	{\hr\hr}mi-	&	{\hr\hr}mi-	&	{\hr\hr}na-\\
Kotos Amarasi	&	{\hr\hr}ʔu-	&	{\hr\hr}mu-	&	{\hr\hr}na-	&	{\hr\hr}ta-	&	{\hr\hr}mi-	&	{\hr\hr}mi-	&	{\hr\hr}na-\\
Molo	&	{\hr\hr}ʔu-	&	{\hr\hr}mu-	&	{\hr\hr}na-	&	{\hr\hr}ta-	&	{\hr\hr}mi-	&	{\hr\hr}mi-	&	{\hr\hr}na-\\
		\lspbottomrule
	\end{tabular}
\end{table}

The internal structure of \tsc{Rote-Meto} provided by a top-down
perspective is mostly straightforward.
There are phonological and morphological innovations which define the two branches.
However, a bottom-up reconstruction provides evidence for a full
series of prenasalised plosives **mb,
**nd, **ŋg, and implosives **ɓ, **ɗ.
While a small number of words with these proto-phonemes have
a PMP etymon, most do not (\citealp[398--402]{Edwards2018a}, \citealp[110--119]{Edwards2023-RMPhon}).\footnote{
		The extent to which
		Proto-Rote-Meto prenasalisation may have correspondences in other
		languages which would allow reconstruction of prenasalisation
		to a higher level remains to be investigated.}
Examples of prenasalised
plosives in \tsc{Rote-Meto} are given in \trf{04-07}.

The changes affecting the Proto-Rote-Meto prenasalised
plosives are shared to different extents between \tsc{Nuclear
Rote} languages and \tsc{Meto} and go against the evidence for \tsc{West Rote-Meto}
provided by the top-down perspective.
Based on these facts \citet{Edwards2018c} proposes that \tsc{Meto}
has shared a period of common development with both west Rote
languages and \tsc{Nuclear Rote} languages -- specifically \tsc{Central
East Rote} (\tsc{Nuclear Rote} excluding Tii and Lole).
Based on the nature of the changes
and the extent to which they have diffused in the lexicons of
the languages involved, \citet{Edwards2018c} proposes that the changes affecting
Dela-Oenale, Dengga,
and \tsc{Meto} occurred once in a common ancestor; Proto-West Rote-Meto.
Pre-Meto then split off
and came into contact with \tsc{Central East Rote} languages.
The changes shared between these groups
then diffused between what were
probably mutually intelligible dialects.

\begin{table}\tcp{\tsc{Rote-Meto} prenasalised plosives}{04-07}
\begin{adjustbox}{width=\textwidth}
	\begin{threeparttable}\st{2pt}
	\begin{tabular}{lllllll}\lsptoprule
local gloss	&	{\hr\hr}betel nut	&	shrimp	&	startle	&	head	&	landslide	&	{\hr\hr}banana\\
PMP	&	{\hr}*buaq	&		&		&		&		&	{\hr}*punti\\
PR-M	&	**\tbr{mb}uah	&	**ka-\tbr{mb}oes	&	**sa\tbr{ŋg}e\tbr{ŋg}er	&	**ka-la\tbr{ŋg}a	&	**\tbr{nd}efa\su{‡}	&	**hu\tbr{nd}i\\ \midrule
Dela-O.	&	{\hr\hr}\tbr{mb}ua	&	{\hr\hr}\tbr{mb}oe-ʔ	&	{\hr\hr}\tbr{ŋg}e\tbr{ŋg}er	&	{\hr\hr}la\tbr{ŋg}a-ʔ	&	{\hr\hr}\tbr{nd}efa	&	{\hr\hr}hu\tbr{nd}i\\
Dengga	&	{\hr\hr}\tbr{mb}ua	&	{\hr\hr}\tbr{mb}oe-ʔ	&	{\hr\hr}\tbr{ŋg}e\tbr{ŋg}e	&	{\hr\hr}la\tbr{ŋg}a-ʔ	&	{\hr\hr}\tbr{nd}efa	&	{\hr\hr}hu\tbr{nd}i\\
Tii	&	{\hr\hr}\tbr{mb}ua	&	{\hr\hr}\tbr{mb}oe-k	&	{\hr\hr}\tbr{ŋg}e\tbr{ŋg}er	&	{\hr\hr}la\tbr{ŋg}a	&	{\hr\hr}\tbr{nd}efa	&	{\hr\hr}hu\tbr{nd}i\\
Lole	&	{\hr\hr}\tbr{mb}ua	&	{\hr\hr}\tbr{mb}oe-k	&	{\hr\hr}\tbr{ŋg}e\tbr{ŋg}e	&	{\hr\hr}la\tbr{ŋg}a	&		&	{\hr\hr}hu\tbr{nd}i\\
Ba{\Q}a	&	{\hr\hr}\tbr{mp}ua	&	{\hr\hr}\tbr{mp}oe-k	&	{\hr\hr}\tbr{ŋg}e\tbr{ŋg}e	&	{\hr\hr}la\tbr{ŋg}a	&	{\hr\hr}\tbr{nd}efa	&	{\hr\hr}hu\tbr{n}i\\
Termanu	&	{\hr\hr}\tbr{p}ua	&	{\hr\hr}\tbr{p}oe-k	&	{\hr\hr}\tbr{ŋg}e\tbr{ŋ}e	&	{\hr\hr}la\tbr{ŋ}a	&	{\hr\hr}\tbr{nd}efa	&	{\hr\hr}hu\tbr{n}i\\
Bokai	&	{\hr\hr}\tbr{p}ua	&	{\hr\hr}\tbr{p}oe-k	&	{\hr\hr}\tbr{ŋ}e\tbr{ŋ}e	&	{\hr\hr}la\tbr{ŋ}a	&	{\hr\hr}\tbr{l}efa	&	{\hr\hr}hu\tbr{n}i\\
Bilbaa	&	{\hr\hr}\tbr{p}ua	&	{\hr\hr}\tbr{p}oe-k	&	{\hr\hr}\tbr{ŋ}e\tbr{ŋ}e	&	{\hr\hr}la\tbr{ŋ}a	&	{\hr\hr}\tbr{l}efa	&	{\hr\hr}hu\tbr{n}i\\
Landu	&	{\hr\hr}\tbr{p}ua	&	{\hr\hr}\tbr{p}oe-ʔ	&	{\hr\hr}\tbr{k}e\tbr{k}er=asa	&	{\hr\hr}la\tbr{k}a-ʔ	&	{\hr\hr}\tbr{nd}efa	&	{\hr\hr}hu\tbr{nd}i\\
Rikou	&	{\hr\hr}\tbr{p}ua	&	{\hr\hr}\tbr{p}oe-ʔ	&	{\hr\hr}\tbr{k}e\tbr{k}e	&	{\hr\hr}la\tbr{k}a	&	{\hr\hr}\tbr{r}efa	&	{\hr\hr}hu\tbr{nd}i\\
Oepao	&	{\hr\hr}\tbr{p}ua	&		&		&	{\hr\hr}la\tbr{k}a	&		&	{\hr\hr}hu\tbr{r}i\\
Roi{\Q}is&	{\hr\hr}\tbr{p}uah	&	{\hr\hr}k\tbr{p}oes	&	{\hr\hr}na-s\tbr{k}e\tbr{k}e	&	{\hr\hr}na\tbr{k}a-f	&	{\hr\hr}\tbr{r}efek	&	{\hr\hh}u\tbr{r}i\\
Kotos&	{\hr\hr}\tbr{p}uah	&	{\hr\hr}\tbr{p}oes,	&	{\hr\hr}na-s\tbr{k}e\tbr{k}e	&	{\hr\hr}ʔna\tbr{k}a-f	&	{\hr\hr}\tbr{k}efan	&	{\hr\hh}u\tbr{k}i\\
	&		&	{\hr\hr}na-k\tbr{p}oe\su{†}	&		&		&		&	\\
Molo	&	{\hr\hr}\tbr{p}uah	&	{\hr\hr}\tbr{p}oes	&	{\hr\hr}na-s\tbr{k}e\tbr{k}e	&	{\hr\hr}ʔna\tbr{k}a-f	&	{\hr\hr}\tbr{k}efaʔ	&	{\hr\hh}u\tbr{k}i\\
Kusa-M.	&	{\hr\hr}\tbr{p}uah, \tbr{p}uha	&	{\hr\hr}\tbr{p}oe-k	&		&	\hr\cg{(boto-f)}	&		&	{\hr\hh}u\tbr{k}i\\
		\lspbottomrule
	\end{tabular}
		\begin{tablenotes}\tnsize
			\item[†]	{Kotos Amarasi \it{na-kpoe} is a verb ‘look/fish for shrimp’.}
			\item[‡]	{Proto-Rote-Meto**ndefa ‘landslide’ is possibly connected
								with PWMP *rəbaq, *rəbas, *Rəbaq ‘collapse, fall down, as a house’.}
		\end{tablenotes}
	\end{threeparttable}
\end{adjustbox}
\end{table}

In the classification in \srf{04-08} at the end of this chapter
we represent only the \tsc{West Rote-Meto} node,
as is appropriate for the family tree model which implies a single common ancestor.
A wave diagram of \tsc{Rote-Meto} is given in \frf{04-02} to
provide a different perspective on the internal relationships of this node.

\begin{figure}[h]\centering
	\includegraphics[width=0.7\columnwidth]{RM-wave.pdf}
	\fcp{\tsc{Rote-Meto} wave diagram}{04-02}
\end{figure}

The changes affecting **mb
and **ŋg provide evidence for further subgroups within \tsc{Nuclear Rote}.
\tsc{Central East Rote} includes all \tsc{Nuclear Rote} languages
except Tii and Lole and is defined by **mb > **mp.
Within this group, all languages except Ba{\Q}a can be placed
in a \tsc{Nuclear Central East Rote} group on the basis of shared
**mp > \it{p} and intervocalic **ŋg > **ŋ.
Within this group, Bilbaa,
Bokai, Landu, Rikou, and Oepao can be placed in a \tsc{Southeast
Rote} group on the basis of word-initial **ŋg > **ŋ.
Landu, Rikou, and Oepao can then be placed together within \tsc{East
Rote} on the basis of subsequent **ŋ > \it{k}.
The changes **mb > \it{mp} > \it{p}
and **ŋg > **ŋ > \it{k} also occur in \tsc{Meto}.

The changes affecting **nd are more complex
and do not align fully with the changes affecting **mb and **ŋg.
The reflexes of Proto-Rote-Meto **nd,
and the environments and languages in which they
occur are summarised in \qf{04-09}.

{\wg\begin{xltabular}{\linewidth}{>{\hspace{-\tabcolsep}}l<{\hspace{-\tabcolsep}}>{\hspace{-\tabcolsep}\enspace}lll>{\raggedright\arraybackslash}X}
\tbx{04-09}	&\mc{4}{l}{Reflexes of Proto-Rote-Meto**nd}	\\
						&a.	&**nd	= \it{nd}					&/{\#}						&Dela-Oenale, Dengga, Tii, Lole, Ba'a, Termanu, Landu\\
						&b.	&**nd	= \it{nd}					&/V{\gap}V				&Dela-Oenale, Dengga, Tii, Lole, Landu, Rikou\\
						&c.	&**nd	> \it{n}					&/V{\gap}V				&Ba'a, Termanu, Korbafo, Bokai, Bilbaa\\
						&d.	&**nd	> \it{r}					&/{\#}						&Rikou\\
						&e.	&**nd	> \it{r} > \it{l}	&/{\#}						&Bokai, Bilbaa\\
						&f.	&**nd	> \it{r}					&/V{\gap}V				&Oepao, Roi'is Amarasi\\
						&g.	&**nd	> \it{r} > \it{k}	&/{\#}, /V{\gap}V	&\tsc{Nuclear Meto} (all \tsc{Meto} except Roi{\Q}is Amarasi)\\
\end{xltabular}\addtocounter{table}{-1}}

Within \tsc{West Rote-Meto}, Dengga
and \tsc{Meto} can be placed together,
as both share **r/*l > **l.
This sound change also affects most Rote languages except Dela-Oenale,
Tii, and \tsc{East Rote} (Landu,
Rikou, Oepao) and it probably spread by diffusion.
Nonetheless, we can propose that it diffused into a common ancestor
of Dengga and \tsc{Meto}.
Three examples of **r > **l in were given in reflexes of *qapəju
‘gallbladder’, *diRus ‘bathe’, and *daləm ‘inside’ in \trf{04-05}.

After \tsc{Dengga-Meto} **r/*l > **l,
\tsc{Meto} then has **l > \it{n} which is the main defining change of this node.
Within Meto, there is a primary split between Roi{\Q}is Amarasi
and \tsc{Nuclear Meto} which contains all other varieties of Meto.
\tsc{Nuclear Meto} is defined by initial **k > \it{k,
ʔ} followed by (**nd >) **r > \it{k}.
Two examples of (**nd >) **r > \it{k} have been seen in reflexes
of **ndefa ‘landslide’ and **hundi ‘banana’ in \trf{04-05}.
Two examples of *k > \it{ʔ}
are given in \trf{04-08}.

Within \tsc{Nuclear Meto}, we can identify a \tsc{Central Meto}
group on the basis of **d > **r > \it{l} which takes place in
all known varieties of \tsc{Nuclear Meto} except Kotos Amarasi, Aam{\Q}abi, and Kusa-Manea.
All varieties of \tsc{Central Meto} except Amanuban
and Molo can then be placed together in a \tsc{North Meto} group
on the basis of **n > \it{l} /lV{\gap}.
Examples of these sound changes are given in \trf{04-08}.

\begin{table}\tcp{\tsc{Meto} sound changes}{04-08}
%\begin{adjustbox}{width=\textwidth}
	\begin{threeparttable}\st{2pt}
	\begin{tabular}{ll@{}l|ll|l@{}l}\lsptoprule
local gloss	&	{\hr}bedbug	&	{\hr}pregnant\su{†}	&	{\hr}ySib	&	{\hr}far	&	{\hr}path, road	&	{\hr\hr}field\\
PMP	&		&	{\hr}*kambu	&	{\hr}*huaji	&	{\hr}*zauq	&	{\hr}*zalan	&	\\
Proto-R-M	&		&	**kambu	&	**waɗi	&	**ka-ɗoo	&	**ɗalan	&	**ɗene\\
Proto-Meto	&	**\tbr{k}etuk	&	**na-\tbr{k}apu	&	**o\tbr{d}iʔ	&	**k\tbr{d}oo	&	**da\tbr{n}an	&	**de\tbr{n}e\\ \midrule
Roi{\Q}is	&	{\hr\hr}\tbr{k}etuʔ	&	{\hr\hr}na-\tbr{k}apuʔ	&	{\hr\hr}o\tbr{r}iʔ	&	{\hr\hr}na-ʔ\tbr{r}oo	&	{\hr\hr}ra\tbr{n}an	&	{\hr\hr}re\tbr{n}e\\
Kotos	&	{\hr\hr}\tbr{ʔ}etuʔ, \tbr{ʔ}etuk	&	{\hr\hr}na-\tbr{ʔ}apuʔ	&	{\hr\hr}o\tbr{r}i-f	&	{\hr\hr}na-ʔ\tbr{r}oo	&	{\hr\hr}ra\tbr{n}an	&	{\hr\hr}re\tbr{n}e\\
Kusa-M.	&	{\hr\hr}\tbr{ʔ}eto	&	{\hr\hr}na-\tbr{ʔ}aup	&	{\hr\hr}o\tbr{r}iʔ	&	{\hr\hr}\tbr{r}oo	&	{\hr\hr}ra\tbr{n}an	&	{\hr\hr}re\tbr{n}e\\
Amanuban	&		&	{\hr\hr}na-\tbr{ʔ}aup	&	{\hr\hr}o\tbr{l}i-f	&	{\hr\hr}ʔ\tbr{l}oo	&	{\hr\hr}la\tbr{n}an	&	{\hr\hr}le\tbr{n}e\\
Amanatun	&	{\hr\hr}\tbr{ʔ}etuʔ	&	{\hr\hr}na-\tbr{ʔ}aap	&	{\hr\hr}o\tbr{l}iʔ	&	{\hr\hr}na-ʔ\tbr{l}oo	&	{\hr\hr}la\tbr{n}an	&	{\hr\hr}le\tbr{n}e\\
Molo	&	{\hr\hr}\tbr{ʔ}etuʔ	&	{\hr\hr}na-\tbr{ʔ}apuʔ	&	{\hr\hr}o\tbr{l}i-f	&	{\hr\hr}ʔ\tbr{l}oo	&	{\hr\hr}la\tbr{l}an	&	{\hr\hr}le\tbr{l}e\\
Amfo{\Q}an	&	{\hr\hr}\tbr{ʔ}etuʔ	&	{\hr\hr}na-\tbr{ʔ}aap	&	{\hr\hr}o\tbr{l}i-f	&	{\hr\hr}na-k\tbr{l}oo-g	&	{\hr\hr}la\tbr{l}an	&	{\hr\hr}le\tbr{l}e-l\\
Timaus	&	{\hr\hr}\tbr{ʔ}etuʔ	&	{\hr\hr}na-\tbr{ʔ}aapw	&	{\hr\hr}o\tbr{l}iʔ	&	{\hr\hr}a-ʔ\tbr{l}oo-gw	&	{\hr\hr}la\tbr{l}an	&	{\hr\hr}le\tbr{l}a-l\\
		\lspbottomrule
	\end{tabular}
		\begin{tablenotes}\tnsize
			\item[†]	{Kotos Amarasi, Amanuban, Molo, and Amfo{\Q}an are
								all also known to have \it{ʔapu-f} ‘womb, uterus’.
								Several reflexes here have final CV → VC metathesis.
								Amanatun and Amfo{\Q}an have metathesis with
								assimilation of the final vowel.}
		\end{tablenotes}
	\end{threeparttable}
%\end{adjustbox}
\end{table}

